%%%%%%%%%%%%%%%%%%%%%%%%%%%%%%%%%%%%%%%%%%%%%%%%%%%
\name{冯云浩}
\contactInfo{13086647572}{Alctrain@163.com}{github/Yunhao-Feng}

\vspace{0.3em}
聚焦 Computer-Use / Coding Agent 的安全、评测与训练，具备 Agent benchmark 构建、运行时监控、tool-use 风险建模、过程级 verifier 设计及基于实机轨迹的 Agentic RL 研究经验。


%%%%%%%%%%%%%%%%%%%%%%%%%%%%%%%%%%%%%%%%%%%%%%%%%%%
\logosection{\faGraduationCap}{教育经历}

\datedline{\textbf{国防科技大学} \quad 博士 \quad 管理科学与工程（工学）}{\dateRange{2024.03}{至今}}
研究方向：Agent 安全、大语言模型安全、鲁棒学习、强化学习优化

\datedline{\textbf{复旦大学} \quad 访问学生}{\dateRange{2025.05}{2025.12}}
研究方向：Agent 安全与大语言模型可靠性评测

\datedline{\textbf{国防科技大学} \quad 硕士 \quad 电子信息}{\dateRange{2022.08}{2024.03}}
研究方向：联邦学习安全、对抗鲁棒性、可信机器学习

\datedline{\textbf{西南财经大学} \quad 本科 \quad 信息管理与信息系统（金融智能与信息管理实验班）}{\dateRange{2018.09}{2022.06}}

%%%%%%%%%%%%%%%%%%%%%%%%%%%%%%%%%%%%%%%%%%%%%%%%%%%
\logosection{\faSuitcase}{实习经历}

\datedline{\textbf{阿里巴巴集团未来生活实验室} \quad 实习算法工程师（Coding Agent 安全方向）}{\dateRange{2025.12}{至今}}
\begin{itemize}[leftmargin=1.5em]
  \item 围绕 Coding Agent 的真实任务执行链路开展安全与可靠性研究，分析 planning、tool use、代码生成与执行过程中的关键攻击面，并参与 Agent 安全评测与攻防验证。
  \item 基于 Agent 实机交互轨迹设计自进化防护模块，探索将运行时风险信号转化为可学习反馈，用于 Agentic RL / post-training 的安全行为优化。
  \item 参与动作前置校验、执行中状态监控与失败轨迹回溯机制设计，支持复杂任务下的过程级监督与防护策略迭代。
\end{itemize}

\datedline{\textbf{京东探索研究院} \quad 算法工程师实习生}{\dateRange{2022.02}{2022.05}}
\begin{itemize}[leftmargin=1.5em]
  \item 参与京东 AutoML 平台表格数据建模模块研发，负责特征预处理、生成与筛选等自动化能力设计与实现。
  \item 交付支持 TB 级结构化数据的自动特征工程框架，并集成至 OmniForce 平台服务多个业务场景。
\end{itemize}

%%%%%%%%%%%%%%%%%%%%%%%%%%%%%%%%%%%%%%%%%%%%%%%%%%%
\logosection{\faBullseye}{岗位相关能力}
\begin{itemize}[leftmargin=1.5em]
  \item \textbf{CUA / Agent 执行环境：} 熟悉 OpenClaw 等 Agent 系统执行流程，具备任务轨迹采集、运行时监控、事件回放与风险分析能力。
  \item \textbf{Benchmark \& Data Construction：} 具备 Agent 安全评测集、复杂长程任务数据与攻击/防护 benchmark 的设计与构建经验。
  \item \textbf{Verifier / Process Supervision：} 具备基于 reasoning trajectory、tool invocation 与 execution outcome 的过程级风险判定与动态拦截经验。
  \item \textbf{Agentic RL / Post-Training：} 正在探索基于真实 Agent 轨迹的自进化安全强化学习，关注 reward shaping、credit assignment 与安全行为优化。
\end{itemize}


%%%%%%%%%%%%%%%%%%%%%%%%%%%%%%%%%%%%%%%%%%%%%%%%%%%

\logosection{\faFlask}{代表性成果}

\textbf{SSE: Safe Self-evolving Reinforcement Learning in Coding Agents} \\
\textit{Under Submission, 2026} \\
提出面向 Coding Agent 的自进化安全强化学习框架，通过 skill extractor 与 agent 协同优化，利用执行轨迹持续提升安全行为表现。

\vspace{0.3em}
\textbf{AgentHazard: A Benchmark for Evaluating Harmful Behavior in Computer-Use Agents} \\
\textit{Under Review, ACM MM 2026 (Dataset Track)} \\
构建 2,653 条 CUA 有害行为评测实例，覆盖多类长程风险任务，用于评估智能体在多步工具调用与上下文累积下的安全中断能力。

\vspace{0.3em}
\textbf{BackdoorAgent: Benchmarking Backdoor Attacks on Agentic Applications} \hfill \textbf{(CCF-A)} \\
\textit{ACL 2026, Findings} \\
提出统一的 Agent backdoor 分析框架，系统研究 trigger 在 planning、memory 与 tool-use 阶段的激活与传播机制。

\vspace{0.3em}
\textbf{SkillTrojan: Backdoor Attacks on Skill-Based Agent Systems} \hfill \textbf{(CCF-A)} \\
\textit{Under Review, 2026, ICML OpenReview Score: 4444 （Weak Accept）} \\
研究 skill-based agent 中的技能级后门攻击，构建 3,000+ backdoored skills 数据集，揭示组合式技能执行中的新型安全风险。


\vspace{0.3em}
\textbf{EZ-ART: Make Agent Red teaming easy} \\
\textit{Under Submission, 2026}

\vspace{0.3em}
\textbf{PaperAsk: A Benchmark for Reliability Evaluation of LLMs in Paper Search and Reading} \hfill \textbf{(CCF-A)} \\
\textit{WWW 2026,\quad (第三作者)}

\vspace{0.3em}
\textbf{FCAT: Federated Causal Adversarial Training} \hfill \textbf{(SCI-1区)}\\
\textit{Knowledge-Based Systems, 2025}

\vspace{0.3em}
\textbf{JAIL: Jailbreak Attacker via Iterative Loss-Guided Beam Search} \hfill \textbf{(SCI-1区)} \\
\textit{Expert Systems with Applications, Under Revision, 2025}


\vspace{0.3em}
\textbf{A Survey of Security Threats in Federated Learning} \hfill \textbf{(SCI-2区)} \\
\textit{Complex \& Intelligent Systems, 2024}






%%%%%%%%%%%%%%%%%%%%%%%%%%%%%%%%%%%%%%%%%%%%%%%%%%%
\logosection{\faCogs}{项目经历}

\datedline{\textbf{XSafeClaw：面向 OpenClaw Agent 的安全运营平台（开源项目）}}{GitHub 开源项目}
\begin{itemize}[leftmargin=1.5em]
  \item 面向 OpenClaw 构建 Agent 运行时监控与风控框架，对 reasoning trajectory、工具调用参数与环境交互行为进行语义级风险识别。
  \item 设计动作前置校验与动态拦截机制，在 tool invocation 前进行风险评估，支持高风险行为的拦截、挂起审批与参数修正。
  \item 构建执行日志回放、事件追踪与可视化分析能力，用于复杂 Agent 行为诊断、安全评测与防护策略迭代。
\end{itemize}

\datedline{\textbf{ClawMonitor: AI Agent Monitor \& Execution Platform}}{GitHub 开源项目}
\begin{itemize}[leftmargin=1.5em]
  \item 开发面向 OpenClaw 的监控与执行平台，支持容器化隔离执行、批量运行、状态监测与轨迹采集，用于 Agent 行为分析和实验复现。
  \item 支持对多轮任务执行链路进行结构化记录，为 benchmark 构建、failure case 分析、verifier 设计及 reward 信号提取提供基础设施。
\end{itemize}

\datedline{\textbf{中文大模型越狱攻击自动生成系统}}{\dateRange{2024.08}{2024.10}}
\datedline{\biInfo{华为杯·第六届中国研究生人工智能创新大赛（大模型安全赛道）全国冠军}{负责人}}{}
设计中文大模型越狱 Prompt 自动生成框架，并完成攻击模板构建与实验验证。

%%%%%%%%%%%%%%%%%%%%%%%%%%%%%%%%%%%%%%%%%%%%%%%%%%%
\logosection{\faCode}{研究方向}
Computer-Use Agents (CUA), Agentic RL / Post-Training, Verifier \& Reward Modeling, Agent Safety \& Red Teaming, Tool-use \& Memory Reliability



%%%%%%%%%%%%%%%%%%%%%%%%%%%%%%%%%%%%%%%%%%%%%%%%%%%
\logosection{\faTrophy}{获奖情况}

\datedline{\textbf{华为杯·第六届中国研究生人工智能创新大赛（大模型安全赛道）全国冠军}}{\dateRange{2024.08}{2024.10}}
\datedline{“兆易创新杯”第十八届中国研究生电子设计竞赛·技术赛三等奖}{\dateRange{2023.05}{2023.08}}
\datedline{“电工杯”全国大学生数学建模竞赛·二等奖}{\dateRange{2021.04}{2021.06}}
\datedline{国防科技大学优秀学员}{\dateRange{2023.12}{2024.12}}
\datedline{国防科技大学国防科技奖学金（校级）}{\dateRange{2023.12}{2024.12}}
